\chapter{Measuring error, uncertainty and the value of redundancy}
\label{ch:evaluation}

% BEGIN SOURCE UNIT M09:00
\section{The experiment must ask the right question}
\label{merge:M09:00}


The proposed ensemble has not yet demonstrated improved interpretation of SFC
circulars. Its design is motivated by inspected literature, the second-circular
exercise and explicit failure analysis. Those are reasons to build and test it,
not evidence that its combination of methods is superior. A credible evaluation
must distinguish a working controller from a useful interpretation method and
from a release process acceptable to the bank.

The main empirical question is whether the method discovers and reports material
interpretation errors and omissions more reliably than simpler alternatives at
comparable cost. A secondary question is whether its reports help reviewers make
better, faster and more traceable decisions. A third concerns resource use:
whether targeted resolution spends effort on questions that can actually be
answered. These questions require different measurements.

The feared failure is a false clean result: a report appears complete and
unambiguous while a material source requirement, alternative reading or factual
assumption has been missed. A system that abstains on every case avoids some false
clean results but may be operationally useless. Evaluation must therefore measure
risk and useful coverage together. Neither raw agreement nor a low average model
uncertainty score is an adequate promotion criterion.

Before a live study, write an evidence contract naming the corpus, baselines,
reference process, primary outcome, uncertainty analysis, hard vetoes, resource
limits and permitted conclusions. Freeze it before examining comparative results.
The manuscript proposes such a protocol; it does not report a completed trial.
Thresholds that depend on the bank's risk tolerance remain decisions for the
responsible owners, not numbers to be invented by the implementation agent.
% END SOURCE UNIT M09:00

% BEGIN SOURCE UNIT P-04-prototype:00
\subsection{The product decision and initial comparison scope}
\label{merge:P-04-prototype:00}

\label{sec:prototype}

The prototype should answer whether a shared executable specification reduces
material mistakes and review effort for a defined class of regulatory changes.
It should also identify the cases that still require judgment and the cost of
maintaining the formalization. A demonstration that translates a clean paragraph
and returns a correct number would establish too little; a project that attempts
all SFC circulars at once would make the source and evaluation problems too large
to diagnose.

The proposed scope is a twelve-week exercise with three connected examples. The
SPI circular and annexes supply the main decision and workflow case. The two
tokenisation circulars supply a version-replacement and actor-specific case.
The fund-marketing circular supplies exceptions and judgment-dependent clauses.
The thematic-review announcement supplies a negative classification example:
the collector must recognize useful regulatory information without inventing an
eligibility rule \citep{sfcspi,sfctoken2023,sfctoken2026,sfcmarketing,sfcreview2025}.
% END SOURCE UNIT P-04-prototype:00

% BEGIN SOURCE UNIT M09:01
\section{Define the unit of evaluation}
\label{merge:M09:01}


A circular contains many questions, and a single question can affect many
transactions. Counting every transaction as an independent interpretation example
would exaggerate sample size when they all depend on the same disputed
classification. The evaluation should distinguish document, provision,
interpretation question, candidate, scenario and transaction levels.

The primary unit for omission detection is a reviewed provision or interpretation
question within a source packet. The source packet includes the circular and
required dependencies as of a declared cut-off. Scenario-level outcomes measure
how candidate rules behave under facts. Document-level summaries measure whether
an apparently complete investigation misses any material issue. These views
complement one another but have different denominators.

For example, the 729 ternary assignments are not 729 independent regulatory
interpretations. They all arise from one selected formula and six abstract inputs.
They are excellent finite-domain engineering coverage for that formula. Treating
them as a sample of 729 circulars would produce meaningless uncertainty intervals
for legal translation accuracy. The 22 named cases likewise share source,
author and modeling assumptions.

The corpus should record these relationships explicitly. Each scenario names its
source packet, provision, interpretation question and reference decision. Each
candidate names its generation method and source version. Analysis can then
aggregate at the appropriate level and use cluster-aware uncertainty estimates
rather than treating related rows as independent observations.

An unresolved reference question is still an evaluation object. The desired
output may be recognition of ambiguity, preservation of supported alternatives
and escalation, rather than selection of a single label. If the dataset forces
such questions into a binary answer, it rewards unjustified certainty and
penalizes the behavior the bank actually needs.
% END SOURCE UNIT M09:01

% BEGIN SOURCE UNIT M09:02
\section{Build the reference before seeing the model's answer}
\label{merge:M09:02}


The reference process begins with independent source inventory. Reviewers enumerate
provisions, footnotes, definitions, exceptions, dates and dependencies. A separate
reader checks the inventory against the retained document. The process records
which materials were unavailable and which judgments depend on them. The corpus
should not silently exclude difficult documents merely because their references
are inconvenient to acquire.

For each selected question, at least two appropriately qualified reviewers form
initial judgments without seeing the candidate system's answer. They state the
reading, assumptions, relevant passages, plausible alternatives and examples that
would distinguish them. This is more work than assigning a label, but it creates
a reference that can test interpretation rather than superficial answer matching.
The exact number and qualifications of reviewers are part of the study plan.

When reviewers disagree, adjudication examines the reasons. A third reviewer or
panel can resolve the issue, retain multiple acceptable readings, or declare the
available authority insufficient. The adjudication record preserves the original
positions and the evidence that changed them. Majority voting among reviewers
should not automatically settle a question of authority or source meaning.

Reviewer independence is documented, not assumed. Shared training, a common case
summary, prior involvement in the implementation and access to model outputs can
all influence judgments. Some collaboration is necessary for adjudication, but
initial blinded judgments should remain distinguishable from the later consensus.
The study reports both stages.

Reference revisions are inevitable when new evidence appears. They require a
versioned correction with a reason. Comparative results are recomputed against
the corrected reference, and the report explains the effect. A correction that
favors one system is not automatically improper, but it must be justified without
using that system's desired score as the reason.

The reference itself is not a universal truth oracle. It is an independently
reviewed account of the declared source perimeter at a time. The study can measure
agreement with that account, sensitivity to its material issues and handling of
its unresolved questions. It cannot establish that no future authority or
previously unknown source could change the interpretation.
% END SOURCE UNIT M09:02

% BEGIN SOURCE UNIT M09:03
\section{Construct a corpus that exposes the expected failures}
\label{merge:M09:03}


A useful corpus contains several kinds of difficulty: explicit Boolean conditions,
exceptions, definitions supplied elsewhere, actor and activity scope, uncertain
classification, amendments, transition periods, event histories, monetary units
and interactions between requirements. The initial public examples can anchor the
engineering harness, but a method comparison needs additional circular families
that were not used to design the prompts or repair logic.

Selection should reflect the intended product perimeter. If the product focuses
on private-bank distribution controls, an evaluation consisting entirely of simple
product-authorization filing instructions would be a weak test of deployment
readiness. Conversely, including arbitrary legal texts outside the supported
language can be useful for measuring abstention, but should not be mixed with
in-scope accuracy without reporting the distinction.

Temporal separation matters. A holdout containing a later amendment can reveal
whether the system improperly reuses an earlier definition or priority. The
source packet must state whether the task is historical interpretation as of a
cut-off or current applicability review. A model with knowledge of a later rule
can otherwise appear to correct a historical answer by importing information
that was unavailable at the assessment date.

Synthetic cases are valuable for controlled defects. They can isolate a missing
exception, an ambiguous attachment or a contradictory fact. Their provenance must
be visible, and their performance must be reported separately from naturally
occurring circular questions. A system can learn the style of seeded mutations
without becoming better at discovering unanticipated omissions.

Include negative controls in which no material ambiguity is supported, as well
as positive controls with known alternative readings. A method that always
invents a dispute should not receive full credit for finding ambiguity. The
reference should distinguish a plausible source-supported alternative from an
arbitrary logical variation. Review burden from unsupported alternatives is an
operational cost worth measuring.

The corpus should also include deliberately incomplete packets. The correct
behavior may be to identify the missing dependency and block a conclusion. These
examples test whether the system recognizes the limits of its evidence. They
must not be scored as ordinary translation errors simply because no final rule
is produced.
% END SOURCE UNIT M09:03

% BEGIN SOURCE UNIT P-05-evidence:00
\subsection{Fifteen cross-circular acceptance scenarios}
\label{merge:P-05-evidence:00}

\label{sec:evidence}

The proposal can be made concrete before connecting to any bank system. The
following cases specify the initial demonstration's intended behavior. Monetary
inputs refer to the adopted financial definitions; the table does not imply that
those facts alone establish SPI status. These are source-based design cases for
independent compliance review. Execution evidence for the selected Java subset
is reported separately below; this broader table is not claimed as completed.

\begingroup\small
\begin{longtable}{P{0.09\textwidth}P{0.41\textwidth}P{0.4\textwidth}}
\caption{Initial acceptance scenarios for the selected scope.}
\label{tab:cases}\\
\toprule Case & Synthetic fact or source change & Required demonstration \\\midrule
\endfirsthead
\toprule Case & Synthetic fact or source change & Required demonstration \\\midrule
\endhead
F1 & Portfolio HK\$40m; qualifying net assets HK\$0. & Establish the financial
condition via the inclusive portfolio comparison.\\
F2 & Portfolio HK\$35m; qualifying net assets HK\$90m. & Establish the financial
condition via the alternative net-asset route.\\
F3 & Portfolio unknown; qualifying net assets HK\$70m. & Leave the financial
condition unresolved; identify missing portfolio evidence.\\
F4 & Portfolio unknown; qualifying net assets HK\$80m. & Establish the financial
condition via the known sufficient route.\\
F5 & Assets HK\$120m, liabilities HK\$15m, primary residence HK\$30m. & Derive
HK\$75m for the illustrated net-asset calculation; assess the separate portfolio
route independently.\\
O1 & HK\$60m joint portfolio with a non-associate; written client allocation 20\%. &
Use HK\$12m for this portfolio-share illustration; explain the agreement and actor roles.\\
K1 & Transaction-experience route established only for one category; a different
category proposed. & Do not propagate a global experience flag; inspect the
alternative qualification routes and selected categories.\\
C1 & A qualifying wealth record, but conservative investment objectives. & Explain
the investment-objective exclusion; wealth alone does not establish SPI treatment.\\
E1 & Designated-account monitoring selected; exposure rises above threshold without
new funds. & Show the relevant FAQ conditions and restrictions; avoid a universal
transaction-block rule.\\
E2 & A withdrawal instruction precedes a proposed streamlined decision in the
adopted event order. & Use the changed consent state and retain the original
instruction and ordering evidence.\\
T1 & A third-party verification clause with an SFC-request trigger. & Preserve the
trigger; distinguish an absent request from missing request records.\\
T2 & Old tokenisation source replaced by the 2026 circular. & Preserve both editions,
record replacement, identify affected rules and reproduce historical decisions.\\
M1 & Product-linked gift versus a discount of fees or charges. & Preserve the stated
exception and the scope of the cited marketing provision.\\
M2 & Daily dealing fund with weekly redemption versus a different daily cut-off. &
Distinguish reduced dealing frequency from administrative procedure.\\
D1 & A thematic-review announcement enters the source collection. & Classify and
track the announcement without inventing a new numerical eligibility criterion.\\
\bottomrule
\end{longtable}
\endgroup

Cases F1--F5 and O1 derive from Annex 1, paragraphs 3.1--3.4. K1 and C1 concern
paragraphs 4--7. E1 requires paragraph 8.3 together with Annex 2, FAQ 6; E2 concerns
the consent provisions and the explicitly adopted event-ordering model. T1--T2
refer to the cited tokenisation editions; M1--M2 to the marketing circular; D1
to the thematic-review announcement
\citep{sfcspiannex1,sfcspiannex2,sfctoken2023,sfctoken2026,sfcmarketing,sfcreview2025}.
This source mapping is part of the acceptance package and should be visible
beside each case in the workbench.
% END SOURCE UNIT P-05-evidence:00

% BEGIN SOURCE UNIT M09:04
\section{Use a baseline ladder with comparable resources}
\label{merge:M09:04}


The simplest baseline is a single structured reading with the same source packet
and a fixed prompt. It establishes what the ensemble adds beyond ordinary
translation. A stronger baseline uses that single reading plus deterministic
schema, coverage and execution checks. This prevents the proposed method from
claiming credit for checks that do not require an ensemble.

A third baseline uses repeated sampling and answer aggregation, inspired by
self-consistency. It receives the same total resource allowance and reports its
aggregation method. A fourth uses diverse initial proposals with a fixed,
deterministic resolution scheduler. The enhanced method adds the proposed
hypothesis-family coverage, targeted questions and optional argumentation or tree
search. This ladder isolates the contribution of each mechanism.

\begin{table}[htbp]
\centering\small
\caption{Comparison arms separate redundancy from search and verification.}
\begin{tabularx}{\textwidth}{P{16mm}YY}
\toprule
Arm & Mechanism & Question it answers\\
\midrule
B0 & One structured proposal & What does ordinary structured translation achieve?\\
B1 & B0 plus deterministic checks & How much comes from source coverage and validation alone?\\
B2 & Repeated proposals with aggregation & Does repeated sampling add useful error detection at equal cost?\\
B3 & Diverse proposals and fixed resolution & Does explicit disagreement handling improve useful coverage?\\
B4 & B3 plus targeted search or argumentation & Does the enhanced scheduler find additional material issues?\\
\bottomrule
\end{tabularx}
\end{table}

Equal resources do not mean merely equal numbers of model calls. Calls differ in
input length, output length, retrieval, latency and cost. The study should record
both token or monetary budget and elapsed time, and present comparisons at
several predeclared budget levels if that is part of the question. Human reviewer
access must also be comparable. Giving only the proposed method a legal reviewer
would confound the method with extra expertise.

Prompt and configuration tuning use a development set. The strongest baseline
should receive reasonable task-specific tuning rather than a deliberately weak
prompt. The test set remains held out until the protocol is fixed. If a test
failure triggers a repair, the repaired method is evaluated on a fresh holdout or
reported as development evidence, not as an untouched final result.

Search methods require their own assumptions audit. A beam width, family quota,
UCT exploration constant or no-progress threshold transferred from puzzles or
another repository is a hypothesis. It can be a starting point, but must not
become the default without target-specific evidence. The deterministic scheduler
is the initial engineering baseline because its behavior is easier to inspect,
not because it has been proven optimal.
% END SOURCE UNIT M09:04

% BEGIN SOURCE UNIT P-04-prototype:04
\subsection{The separate comparison of whole review workflows}

The baseline ladder above compares automatic interpretation methods. A second
experiment is needed for the product decision: does the complete workbench help
people carry a change through review? The four workflow conditions below answer
that broader question. They must not be pooled with repeated model calls as if
both experiments had the same observational unit.


\label{merge:P-04-prototype:04}


The principal engineering question is whether the workbench produces a more
reviewable and dependable implementation of the selected requirements than the
current process. The baseline must therefore be a faithful version of that
process, with the same documents, permitted reference material, business scope and
definition of a completed change.

Use four conditions: the existing manual process; a source-grounded LLM
summary/checklist; a hand-authored specification with the workbench; and the same
workbench with model-assisted drafting. The summary baseline tests whether a
simpler product achieves the benefit. The two formalized conditions isolate the
incremental value of the model. Optional proof/search enhancements form a later
arm under the same information and budget constraints. Preparation of formal
rules, examples and reusable definitions must be counted, rather than treated as
free information available only to the assisted condition. This addresses the
information-budget distinction exposed by the tax-reasoning literature
\citep{jurayj2026}.

Independent compliance reviewers should derive expected outcomes and required
actions from the selected sources before seeing the candidate implementation.
Where interpretation is contested, the expected outcome can be a review request
or an explicitly adjudicated reading. Ambiguity must not be removed merely to
obtain a convenient deterministic label. Cases generated from the formal model
serve a different role: they exercise its paths and mutations, and should be
reported separately from the independent source-derived cases.

Material errors include applying a rule to the wrong actor or category, allowing
a scoped action after a relevant withdrawal, losing a required obligation,
inventing a regulatory restriction, and converting decision-relevant missing
information into an established fact. The exact taxonomy and severity must be
agreed in weeks 1--2. A correct financial subcondition must not be scored as an
incorrect trade permission merely because the test harness confuses those outputs.

\begingroup\small
\begin{longtable}{P{0.25\textwidth}P{0.39\textwidth}P{0.26\textwidth}}
\caption{Proposed evidence contract for the prototype evaluation. These are
prospective criteria, not results.}\label{tab:evaluation}\\
\toprule Evidence & Definition and use & Decision role \\\midrule
\endfirsthead
\toprule Evidence & Definition and use & Decision role \\\midrule
\endhead
Material interpretation and decision defects & Independently adjudicated errors
in source coverage, rule meaning, resulting decisions or obligations. No unresolved
critical defect may be carried into a proposed operational scope. & Primary quality
criterion; blocks promotion of the affected scope.\\
Review effort & End-to-end person-time for a completed, accepted change, including
source work, corrections, escalations and reusable-rule preparation. & Primary
business criterion, conditional on quality.\\
Unknown and review outcomes & Frequency, reason, downstream handling and time to
resolution; distinguish necessary abstention from unhelpful failure. & Explanatory
diagnostic and business cost; incorrect silent resolution is a quality defect.\\
Source and dependency completeness & Every in-scope rule has inspected support and
resolved imports; every required selected clause is accounted for. & Release veto
for the selected corpus. Does not establish completeness of all HK regulation.\\
Adapter agreement and mutations & Supported semantics agree; targeted changes in
operators, scope, dates and exceptions are detected by the required checks. &
Engineering veto on the adapter; mutation scores alone do not establish legal fidelity.\\
Reproducibility & A saved rule version, fact snapshot and event sequence reproduce
the original result and explanation references. & Engineering veto on historical replay.\\
Latency and token/runtime cost & Measure on the specified local workload, including
failed drafts and review retries. & Explanatory and operating-budget evidence; no
quality trade-off is implied.\\
\bottomrule
\end{longtable}
\endgroup

W10 proposes at least 30 independently adjudicated change tasks, provisionally
10 for development and 20 held out. This is a feasibility budget assumption,
not a statistical power calculation. Split by circular family and amendment
chain; if the five initial documents cannot supply a credible separation, expand
the corpus before making a generalization claim. Two compliance readers derive
the expected scope, duties and case outcomes; contested interpretations remain
a separate stratum rather than disappearing from the denominator.

Report paired error and reviewer-time differences, with intervals that account
for reviewer and document effects. Repeated calls on one paragraph are not
independent regulatory changes. Assign comparable tasks, rotate conditions where
feasible and record prior exposure to avoid giving the assisted condition an
unearned learning advantage.

Zero unresolved critical false approvals or missed duties is a proposed release
screen on the adjudicated cases. Passing it does not estimate a zero population
error rate. Criticality is defined with the bank before unblinding. If too few
independent families remain for a defensible comparison, report descriptive
results and the additional sample needed rather than rank viable approaches.

A small pilot may establish useful failure cases without supporting a precise
speed ranking. Zero observed material errors in a finite synthetic set does not
establish a zero error rate in the bank's workload. Similarly, a high mutation
score says the selected perturbations were detected; it does not show that every
missing legal requirement would be found. Report these limits directly and use
the evidence to select the next scope of validation.

The evaluation should pause if its source inventory is wrong, its expected
outcomes are not independently supportable, records are corrupted, or conditions
receive materially different information without accounting for it. Those are
problems with the evaluation. A failed drafting strategy or event adapter is a
candidate failure: it should trigger a specified repair and repeat of the affected
checks, rather than abandonment of the whole product direction.
% END SOURCE UNIT P-04-prototype:04

% BEGIN SOURCE UNIT M09:05
\section{Measure omissions, alternatives and useful abstention}
\label{merge:M09:05}


Let $Q$ be the set of independently reviewed material questions in the evaluation
corpus. For each $q\in Q$, the reference records required conditions, accepted or
plausible readings and whether available evidence resolves the question. A
material omission occurs when the system's report fails to represent a required
condition or unresolved question within its claimed scope. This definition
includes missing uncertainty, not only incorrect final formulas.

A false clean result occurs when the system reports a question as resolved and
eligible for review without disclosing a material unresolved defect identified by
the reference. The denominator is the number of questions or documents declared
clean, depending on the reported metric. A system that declares fewer clean cases
may reduce this rate while increasing manual workload. The report must show both.

Define useful coverage as the fraction of eligible questions for which the system
produces a reference-supported, sufficiently complete result without unresolved
material defects. Define selective risk as the error rate among the questions it
chooses to resolve. A risk-coverage curve varies the decision policy and shows the
tradeoff. The policy must not be tuned on the final test labels after the fact.

Alternative-reading recall measures whether the report contains the reference's
material supported alternatives. It should be assessed by meaning and consequence,
not string matching. Alternative precision measures how many proposed alternatives
are judged relevant and supported. These measures expose a method that increases
recall by overwhelming reviewers with arbitrary possibilities.

Resolution quality has a different denominator: detected discrepancies. Measure
whether the selected action obtains the needed evidence, produces a valid repair,
correctly escalates, or falsely closes the issue. A successful HTTP response is
not a successful resolution. The action's declared question and expected evidence
type make this classification possible.

Report completeness is largely an engineering metric. It checks that every issue,
role, candidate disposition and resource outcome appears. It can be tested
exhaustively over bounded controller fixtures. It does not establish that the
issues themselves cover every material legal question. Keeping these metrics
separate prevents a perfect report-rendering score from being presented as
perfect interpretation accuracy.
% END SOURCE UNIT M09:05

% BEGIN SOURCE UNIT M09:06
\section{Severity and direction of error need explicit treatment}
\label{merge:M09:06}


Not every error has the same consequence. Missing a prohibition may permit an
impermissible offer; inventing a prohibition may block a legitimate activity or
mislead the bank about its obligations. Both are errors relative to the accepted
interpretation, even if one appears operationally conservative. The evaluation
should report their directions rather than collapsing them into a single accuracy
number.

Severity categories must be defined before scoring. They can consider whether the
error changes an action, affects many clients, is hard to detect downstream or
undermines the source and approval chain. Assigning weights is a policy choice
requiring the bank's input. The study can report unweighted counts and separate
severity strata until those weights are approved.

A weighted average can hide a rare severe error among many trivial successes.
Therefore severe false clean outcomes should also be reported as individual cases
with their mechanisms. A predeclared severe-error veto can block promotion even
when average performance is favorable. That veto should not automatically stop
all research: if the failure reveals a repairable omission detector, the next
planned repair phase may remain justified.

The report should distinguish an invalid experiment from a failed candidate.
Leakage of test labels into prompts invalidates the comparison and may require
rebuilding the evaluation. A candidate missing an exception on a valid heldout
case is evidence against that candidate's readiness. An unavailable provider is
an operational failure that may prevent estimating interpretation performance.
These outcomes should not all be summarized as ``the experiment failed.''

Reviewer workload also has severity. A report containing hundreds of unsupported
alternatives may conceal the important issue through overload. Measure review
time, number of evidence inspections, unresolved-question quality and corrections
made by reviewers. A method that catches more errors but makes the review process
unmanageable needs a different presentation or search policy, not an unqualified
claim of superiority.
% END SOURCE UNIT M09:06

% BEGIN SOURCE UNIT M09:07
\section{Zero observed errors does not mean zero risk}
\label{merge:M09:07}


Suppose a study observes no errors in $n$ independent, identically distributed
trials with unknown error probability $p$. Under this simplified binomial model,
the probability of seeing zero errors is $(1-p)^n$. A one-sided upper confidence
bound at level $1-\alpha$ solves $(1-p)^n=\alpha$, giving
\begin{equation}
 p_{\rm upper}=1-\alpha^{1/n}.
 \label{eq:zero-errors}
\end{equation}
For a 95 percent bound and $n=100$, this is approximately $2.95$ percent. With
$n=1000$, it is approximately $0.30$ percent. These calculations show why a small
perfect test set cannot establish a very small error rate.

The independence assumption is substantial. One hundred scenarios derived from
one circular do not generally behave like one hundred independent circulars.
Shared wording, definitions and prompts create clusters. The calculation is
therefore illustrative unless the trial definition and sampling design justify
the model. It should not be applied to the 729 truth-table assignments as a legal
error bound.

For method comparisons, paired evaluation on the same source questions reduces
irrelevant variation. A question-level difference records whether one method
finds a material issue that the other misses. Confidence intervals or permutation
procedures should preserve the document or question clustering appropriate to the
sampling design. Resampling individual rows while ignoring their shared circular
can produce overly narrow intervals.

Multiple stochastic generations require multiple seeds or repeated provider runs,
with the seed and configuration recorded where supported. Differences from one
run are descriptive. If a model API does not guarantee reproducibility from a
seed, the study records independent request repetitions and their limitations.
The retained responses permit audit even when exact regeneration is impossible.

Tail measures such as maximum review time or the 99th percentile of action cost
need particularly cautious interpretation. A few difficult circulars can dominate
them, and small samples give unstable estimates. They are useful diagnostics and
capacity-planning signals, but a superiority claim requires a predeclared analysis
with sufficient observations and uncertainty evidence.

The evaluation report should state which candidates failed hard safety or
integrity screens, which remain viable, whether any ranking is statistically
supported, and which differences are descriptive only. Passing a hard screen is
not evidence that one viable method is better than another. This discipline is
especially important when the apparent gain comes from a small number of rare
severe errors.
% END SOURCE UNIT M09:07

% BEGIN SOURCE UNIT M09:08
\section{Calibration asks a narrower question than truth}
\label{merge:M09:08}


A calibrated score has a defined relationship to an observed event in a target
population. For example, among cases assigned a particular probability range,
the fraction agreeing with the independently reviewed reference should be
consistent with that range, within sampling uncertainty. Calibration is not a
property conferred by using the word confidence in a JSON field.

The target event must be precise. ``The interpretation is correct'' is too broad
unless the reference process, source perimeter and acceptable alternatives are
specified. A more measurable event might be ``the proposed paragraph-level rule
has no material discrepancy against the adjudicated reference for this corpus.''
Even that event depends on the quality and scope of the reference.

Calibration data must be separate from the data used to tune the score and from
the final evaluation set. A score trained on one circular family may be
miscalibrated on another, especially when definitions, authority structure or
source quality change. Model and prompt updates can also invalidate an earlier
calibration. The calibration record therefore binds the method and population
versions.

Reliability diagrams and proper scoring rules can diagnose probability estimates,
but a visually plausible diagram with few cases is weak evidence. Calibration
also does not ensure discrimination: a model that assigns the base rate to every
case can be calibrated yet unhelpful for selecting review priorities. The study
should assess both risk ranking and calibration, with uncertainty intervals.

Semantic entropy can nominate uncertain answers, as discussed in
Chapter~\ref{ch:search}. Its published evaluations use particular datasets and
correctness proxies \citep{semantic2023}. A high area under a ranking curve in
those tasks does not establish calibrated legal error probabilities. In the
workbench, entropy may help prioritize review after a target-specific study, while
mandatory source and release checks remain categorical requirements.

Low entropy is particularly dangerous under common-mode omission. Every sampled
answer can concentrate on the same incomplete reading. The independent inventory
and source review remain necessary even if an uncertainty score appears excellent
on average. A useful evaluation explicitly includes these unanimous-error cases
rather than assuming that disagreement accompanies every mistake.
% END SOURCE UNIT M09:08

% BEGIN SOURCE UNIT M09:09
\section{Conformal prediction can support sets under stated assumptions}
\label{merge:M09:09}


Conformal prediction offers a principled way to construct prediction sets from
calibration examples \citep[section 1.1 and appendix D]{conformal2022}. It is
relevant because the proposed product should sometimes return several candidates
rather than a forced answer. Its guarantee, however, concerns a specified
statistical prediction problem under exchangeability. It is not a theorem about
all possible interpretations of English.

To make the mechanism concrete, suppose each task has a finite reviewed label
space $\mathcal Y$, such as a small set of predetermined interpretation classes.
A fitted scoring function $s(x,y)$ assigns larger values when label $y$ looks less
compatible with task $x$. The scoring function is fixed before calibration.
For $n$ calibration tasks with reviewed labels $(X_i,Y_i)$, compute
$S_i=s(X_i,Y_i)$.

Choose a target miscoverage level $\alpha$ and define
\begin{equation}
 k=\left\lceil(n+1)(1-\alpha)\right\rceil,
 \qquad
 \widehat q=\begin{cases}
 S_{(k)},&k\leq n,\\
 +\infty,&k=n+1,
 \end{cases}
 \label{eq:conformal-quantile}
\end{equation}
where $S_{(k)}$ is the $k$th smallest calibration score. For a new task $x$, return
\begin{equation}
 \Gamma(x)=\{y\in\mathcal Y:s(x,y)\leq\widehat q\}.
 \label{eq:conformal-set}
\end{equation}
The finite-sample correction uses $n+1$, not simply the empirical $1-\alpha$
quantile with an unspecified interpolation rule. The inclusive inequality also
matters when scores tie.

Why does this work? Conditional on the separately fitted scoring rule, if the
$n$ calibration examples and the new example are exchangeable, the position of
the new true-label score among the $n+1$ scores is symmetric. In the no-tie case,
its rank is uniform. Excluding it requires a rank beyond $k$, whose probability
is at most $\alpha$ by the definition of $k$. With ties and inclusive sets, the
basic coverage statement is conservative. This gives marginal coverage over a
new exchangeable example and the calibration sample.

For $n=19$ and $\alpha=0.1$, $k=18$, so the threshold is the eighteenth ordered
calibration score. For $n=4$ and $\alpha=0.1$, $k=5$, requiring the infinite
threshold and potentially the entire label space. The latter example makes a
practical point: a very small calibration set cannot support a narrow set with a
strong finite-sample coverage target through this construction.

The guarantee does not say that each particular circular has a 90 percent chance
of containing the right reading in its returned set. Nor does it guarantee
coverage conditional on every rare ambiguity type. Marginal coverage can coexist
with poor performance on a small important subgroup. The distinction between
marginal and conditional coverage is discussed explicitly in the tutorial
\citep[section 3]{conformal2022} and is central to a high-stakes application.

The label space is another limitation. If the generator has never proposed a
material reading and the label space contains only generated candidates, a set
selection method cannot include that missing reading. One can add an
``unrepresented interpretation'' or escalation outcome to a reviewed task
formulation, but that changes the prediction problem and requires reference data.
Conformal calibration does not establish semantic completeness of open-ended
candidate generation.

The appropriate early experiment is therefore narrow: use a fixed, reviewed
classification task or a defined error-detection score, calibrate on separate
source families where exchangeability is defensible, and measure set size,
coverage and subgroup behavior. Source amendments, provider changes and shifts in
document type should trigger renewed assessment. The inspected official tutorial
code is useful for checking quantile and set construction, but its image examples
are not a ready-made legal calibration protocol.
% END SOURCE UNIT M09:09

% BEGIN SOURCE UNIT M09:10
\section{Search coverage and semantic completeness are different}
\label{merge:M09:10}


A controller can report that it explored every node in a finite generated graph.
That is a meaningful computational fact. It does not establish that the graph
contains every supported legal interpretation. The generator, inventory and
construction operators determine what enters the graph. Evaluation must examine
both the search procedure and the generation boundary.

For each reference alternative, record whether it was generated, retained,
investigated, correctly dispositioned and represented in the final report. This
stage-by-stage accounting locates the failure. If an alternative was never
generated, changing the scheduler may not help. If it was generated but pruned by
a global beam before investigation, family preservation may be useful. If it was
investigated but falsely defeated by a weak argument, the evaluation criterion
needs repair.

The same accounting applies to omissions. A source unit can be missed at
acquisition, extraction, inventory, normative classification or rule mapping.
Treating all such failures as model reasoning errors hides the actionable cause.
The ensemble's modular design makes it possible to test which redundant activity
actually contributed detection.

Ablations remove one component at a time under comparable budgets: the independent
inventory, the alternative reader, the argument critic, the witness generator or
the no-progress scheduler. If removing a component changes both resources and
information access, the interpretation of the ablation is limited. The study
should state whether it measures the component's total contribution or its value
at fixed cost.

Search completeness claims must name the finite object. ``All twenty-four retained
candidates were compared on the declared Boolean domain'' is testable. ``All legal
readings were considered'' is unsupported. The report should also expose the
unexplored frontier, discarded unsupported proposals and resource-limited
comparisons so that reviewers can judge whether more investigation is justified.
% END SOURCE UNIT M09:10

% BEGIN SOURCE UNIT M09:11
\section{Evaluate the reviewer, not only the model}
\label{merge:M09:11}


The final product is a workbench used by people. A study should therefore compare
review with and without the ensemble dossier, using a design that controls for
case difficulty and learning effects. Reviewers should not see the same case in
both conditions in an order that makes the second condition trivially easier.
A counterbalanced or appropriately randomized assignment can reduce this problem.

Outcomes include detection of material defects, quality of the final rationale,
time to a defensible decision, number of unnecessary escalations and ability to
identify remaining uncertainty. Faster review is valuable only if the substantive
quality is preserved. A dashboard that encourages reviewers to accept a misleading
clean badge can reduce time while increasing risk.

The study should inspect automation bias directly. Include cases where the
leading model answer is wrong, where a minority candidate is right relative to
the reference, and where no single answer is justified. Observe whether reviewers
inspect the opposing argument and source. The purpose is not to trick staff, but
to learn whether the interface supports the intended skeptical review.

Reviewer feedback also tests self-containment. Can a compliance reader understand
the formal rule through its controlled-language rendering and examples? Can an
engineer identify the required Java behavior without consulting an external paper?
Can both explain why a solver result does not resolve an uncertain definition?
These comprehension tasks are more informative than asking whether the report
looks professional.

Any study involving actual bank staff, confidential materials or production
workflows requires the bank's normal authorization and data controls. The public
prototype can first use synthetic tasks and qualified independent readers.
Its findings should be reported at that level. A successful public-source
usability exercise is not evidence that every private-bank process has been
validated.
% END SOURCE UNIT M09:11

% BEGIN SOURCE UNIT P-04-prototype:06
\section{Economic value and operating cost}
\label{merge:P-04-prototype:06}


The business case should use measured labor and maintenance effort. Let $K$ be
the initial cost of constructing the selected corpus, workbench and integration.
Let $M$ be the average manual cost of a comparable completed change, and let $A$
be its assisted cost including human review, model use, operation and allocated
maintenance. If those costs remain comparable across changes and $M>A$, the
simple break-even count is
\begin{equation}
 n_{\mathrm{break\text{-}even}}=\frac{K}{M-A}.
 \label{eq:cost}
\end{equation}
This follows by comparing manual cost $nM$ with assisted cost $K+nA$ and solving
$nM=K+nA$. It is a planning relation, not an estimated return. If $M\leq A$, this
labor-saving model alone supplies no break-even case. Material-error prevention,
audit response and amendment speed may have additional value, but they should be
measured separately rather than assigned unsupported monetary benefits.

Reusable definitions can lower later costs, while new product classes and
ambiguous amendments can raise them. The prototype should record both effects.
Paper benchmark token costs and tax penalties are not substitutes for the bank's
staffing costs, risk definitions or change volume. A credible investment decision
may be to continue only for a narrow class of rule changes where the measured
benefit is clear.
% END SOURCE UNIT P-04-prototype:06

% BEGIN SOURCE UNIT M09:12
\section{Promotion decisions and the next evidence needed}
\label{merge:M09:12}


The study should predeclare a small set of promotion criteria and hard vetoes.
A severe false clean result, broken source identity, uncounted external action,
unauthorized release or incomplete mandatory report can veto promotion. Average
runtime, model score and cost are explanatory until the correctness and integrity
conditions are met. A method cannot trade away a required release invariant for
better throughput.

Among methods that pass the hard screens, ranking requires the predeclared
uncertainty analysis. If the data do not distinguish them, report them as viable
under current evidence. A descriptively favorable method may justify a longer
study or an optional feature, but not a new default merely because its mean score
is highest in a small sample.

\begin{table}[htbp]
\centering\small
\caption{Decision record required after a substantive method comparison.}
\begin{tabularx}{\textwidth}{P{39mm}Y}
\toprule
Decision field & Required interpretation\\\midrule
Primary criterion & State the observed result and uncertainty against the frozen criterion.\\
Veto status & Identify severe errors, integrity failures and invalid comparisons before discussing speed.\\
Viable candidates & List methods that remain eligible for further evaluation; do not imply an unsupported ranking.\\
Main uncertainty & Distinguish reference ambiguity, sampling error, distribution shift and untested implementation.\\
Next justified action & Promote a bounded use, extend the study, repair a mechanism, or rebuild an invalid harness.\\
Limit of conclusion & State the source perimeter, task population and deployment claims not established.\\
\bottomrule
\end{tabularx}
\end{table}

A failed candidate should trigger the repair that its failure supports. Missing
source units suggest improving inventory and extraction diversity. Losing a
minority interpretation suggests changing family retention. Repeated ineffective
actions suggest a better question-selection policy. A contaminated reference
requires repairing the evaluation before comparing methods again. These are
specific next steps, not reasons to declare the entire research direction either
successful or impossible.

Every serious run should retain a manifest with source and code identities,
commands, environment, model/provider versions, configurations, seeds or request
repetitions, wall time, resource use, plan and result paths. The result note should
include the strongest alternative explanation and the evidence that would overturn
its conclusion. This makes the evaluation reproducible enough to guide product
decisions rather than merely produce an attractive benchmark table.

The immediate conclusion remains limited: the proposed method is ready to be
implemented as a bounded, testable extension, and its empirical legal-interpretation
performance remains to be established. The next chapter describes how the bank
can preserve that distinction through shadow operation, source changes and
ongoing review.
% END SOURCE UNIT M09:12

